\section{Formalization plan and dependency separation}
\label{sec:formalization}

This appendix is a plan, not a collection of machine-checked declarations.
No uncompiled Lean file with placeholder proof terms is presented as a
proof artifact.

\subsection{Finite cubic theorem}

The main algebraic theorem can be formalized independently of analytic
$L$-functions. Start with the finite type $X=\F_3^r$, its dual $V$, and
a nonzero coefficient function $c:X\to\Z$ with zero sum. Define its
positive mass and the transform in the Eisenstein integers. A natural
dependency order is:
\begin{enumerate}
\item Finite Fourier orthogonality and inversion; support symmetry and
the zero value at the identity.
\item The zero-sum triple lemma, including the exact odd/even
$1-\zeta_3$ valuations and their mass bounds.
\item Coordinate interpolation on $\{0,1,2\}$ over the Eisenstein ring
with 2 inverted; coefficient divisibility and reduction to a low-degree
ternary polynomial.
\item The needed minimum and second-weight Reed--Muller theorems,
including the affine support-containment statements. Their mathematical
citations are explicit, but availability as matching Lean declarations
has not been established in this session.
\item The piecewise lower bound, exact coset constructions, and all
equality cases via Fourier inversion.
\end{enumerate}
Steps 1--5 yield a purely finite theorem. Its Dirichlet interpretation
then needs surjectivity of the residue-class evaluation map and the prime
number theorem in arithmetic progressions. The finite-prime statements
separately need the smoothed explicit-formula interface and the exported
Dirichlet GRH. The imported GRH does not automatically formalize those
intermediate estimates.

\subsection{Branch recovery}

The second development can also be kept independent of GRH. Its core is
formal substitution and division in $F[[X,Y]]$, the discrete valuation on
$F[[X]]$, a finite weighted greedy ordering, and an exact Newton basis.
After the kernel identity, filtered linear algebra gives the target-jet
dimension and recovery criterion. Lagrange interpolation supplies the
$d$-branch certificate and sharp recording budgets.

The error/erasure theorem should be formalized for the quotient target
explicitly. Two different source series may have the same target jet,
so ordinary injectivity of the whole observation map is the wrong
specification. One can define target distance as the smallest support of
an observation word arising from a series outside $(X,Y)^d$, then prove
the distance criterion and the tree recursion.

\subsection{What the supplied Lean theorem does and does not discharge}

After the finite and analytic interfaces above have been formalized,
the supplied declaration
\begin{center}
\code{LFunction\_ne\_zero\_of\_half\_lt\_re}
\end{center}
or the critical-strip theorem can discharge the Dirichlet-GRH hypothesis of
the counting results. It is unnecessary for the exact extremal density
or branch-kernel theorems. Neither the absence of forbidden tokens in
the vendor source nor the success of the finite Python checks should be
used as a substitute for this future formal proof.
